% Source: https://artofproblemsolving.com/community/c3h2304469_distinguishability_help
% Stable identity: aops-topic-2304469-q01
\textbf{(a)} Each ball independently chooses one of $3$ boxes: $3^5=\boxed{243}$.

\textbf{(b)} We need the solutions of $a+b+c=5$ in nonnegative integers; by stars and bars there are $\dbinom{7}{2}=\boxed{21}$.

\textbf{(c)} Now only the grouping of the balls matters: we split the set of $5$ balls into at most $3$ nonempty groups.
\begin{itemize}
\item $1$ group: $1$ way.
\item $2$ groups: each of the $2^5-2=30$ ways to colour the balls with two colours, both used, gives a split, and each split arises twice: $15$ ways.
\item $3$ groups: sizes $3,1,1$: choose the group of $3$, $\dbinom53=10$ ways; sizes $2,2,1$: choose the single ball ($5$ ways) and split the other $4$ into two pairs ($3$ ways), $15$ ways. Total $25$.
\end{itemize}
In all, $1+15+25=\boxed{41}$.

\textbf{(d)} Only the multiset of group sizes matters, i.e. the ways to write $5$ as a sum of at most $3$ positive parts: $5$, $4+1$, $3+2$, $3+1+1$, $2+2+1$. There are $\boxed{5}$.
